% ---------------------------------------------------------------------
%  06 Results
%  Every subsection: state the claim as the heading sentence, then the
%  table/figure that supports it, then one sentence of what it means.
% ---------------------------------------------------------------------
\section{Results}
\label{sec:results}

\subsection{2D perception}
\label{subsec:results_2d}

\begin{table}[t]
  \centering
  \caption{\Todo{Ball detection accuracy. Columns: method, backbone, temporal
  context $T$, localisation error (px), detection rate.}}
  \label{tab:ball_detection}
  % \input{assets/tables/ball_detection.tex}
  \begin{tabular}{lccc}
    \toprule
    Method & $T$ & Err. (px) $\downarrow$ & Det. rate $\uparrow$ \\
    \midrule
    \Todo{fill} & -- & -- & -- \\
    \bottomrule
  \end{tabular}
\end{table}

\begin{table}[t]
  \centering
  \caption{\Todo{Court keypoint detection and homography quality.}}
  \label{tab:court_detection}
  \begin{tabular}{lcc}
    \toprule
    Method & PCK $\uparrow$ & Reproj. err. (px) $\downarrow$ \\
    \midrule
    \Todo{fill} & -- & -- \\
    \bottomrule
  \end{tabular}
\end{table}

\subsection{3D ball trajectory (\BLCS{})}
\label{subsec:results_blcs}
\Todo{Compare against triangulation baseline. Report error vs. number of
views in which the ball is detected --- that curve is the interesting result.}

\subsection{Player root position and orientation (\PLCS{})}
\label{subsec:results_plcs}
\Todo{Position error (m) and rotation error (deg). Report the rotation error
distribution, not just the mean: the failure mode of interest is bimodal
($180^\circ$ flips), and a mean hides it.}

\begin{figure}[t]
  \centering
  % \includegraphics[width=\linewidth]{assets/figures/plcs_rotation_hist.pdf}
  \fbox{\parbox[c][0.22\textheight][c]{\linewidth}{\centering
    \Todo{plcs\_rotation\_hist.pdf --- histogram of rotation error, before vs.
    after the angle loss}}}
  \caption{\Todo{Caption naming the bimodality explicitly.}}
  \label{fig:plcs_rotation_hist}
\end{figure}

\subsection{End-to-end court-frame reconstruction}
\label{subsec:results_e2e}
\Todo{The headline result: ball and articulated players in one metric frame.
Include qualitative reconstructions and at least one honest failure case.}

\subsection{Ablations}
\label{subsec:results_ablation}

\begin{table}[t]
  \centering
  \caption{\Todo{Ablation table. One knob per row, everything else held
  fixed; say what is held fixed in the caption.}}
  \label{tab:ablation}
  \begin{tabular}{lcc}
    \toprule
    Variant & Pos. err. (m) $\downarrow$ & Rot. err. (deg) $\downarrow$ \\
    \midrule
    Full model      & -- & -- \\
    \Todo{$-$ knob} & -- & -- \\
    \bottomrule
  \end{tabular}
\end{table}

\subsection{Sim-to-real transfer}
\label{subsec:results_sim2real}
\Todo{Train on simulation, evaluate on real clips. Quantify the gap and
attribute it: input-distribution mismatch (court keypoint count/noise) versus
motion-distribution mismatch.}
